% language=us

\environment context-2025-style

% \setuptyping
%   [style=\switchtobodyfont[7pt,tt]]

\definelayer
  [extratext]
  [preset=rightbottom,
   height=\textheight,
   width=\textwidth]

\setupbackgrounds
  [text]
  [background=extratext]

% \definehighlight[notabene][style=bf]
\definehighlight[notabene][style=\glyphweight100\underbar]

\setupalign[flushleft]

\setuptyping
  [align=hangright,
   keeptogether=paragraph,
   numbering=]

\startdocument
  [title={compact fonts},
   banner={a practical approach},
   location={context\enspace{\bf 2025}\enspace meeting}]

\starttitle[title=Eight bit(maps)]

\startitemize
\startitem
    Originally \TEX\ only supported bitmap fonts. These were generated
    on demand at the right resolution, determined by target device and
    scale.
\stopitem
\startitem
    When \TYPEONE\ fonts showed up \POSTSCRIPT\ backend drivers made
    sure that we would use those.
\stopitem
\startitem
    The limited set of original fonts came with design sizes so there had to be
    some mapping (and choice) to desired size.
\stopitem
\startitem
    Those \POSTSCRIPT\ fonts seldom had design sizes translated into multiple
    instances. Scaling was easy.
\stopitem
\startitem
    In both cases there had to be suitable metric \TFM\ files, possibly
    accompanied by virtual font definitions.
\stopitem
\startitem
    Fonts can be extended (wider), slanted (angle) and subtle horizontal
    scaling can emulate hz. This is also map file driven (backend).
\stopitem
\stopitemize

\stoptitle

\starttitle[title=Going wide]

\startitemize
\startitem
    When wide fonts came available \PDFTEX\ was able to use them but still within
    eight bit subsets.
\stopitem
\startitem
    Both \XETEX\ and \LUATEX\ made sure that we could handle any \OPENTYPE\ font
    in \UNICODE\ aware mode.
\stopitem
\startitem
    The metrics and other properties come from the font itself. There are very
    few design sized fonts, although one can discuss if variable fonts suit that
    description (quality wise).
\stopitem
\startitem
    But, \XETEX\ and \LUATEX\ still use instances in the engine, pre-calculated
    dimension vectors.
\stopitem
\startitem
    The backends have to support extend, slant and hz.
\stopitem
\startitem
    But we can do better.
\stopitem
\stopitemize

\stoptitle

\starttitle[title=Math]

\startitemize
\startitem
    In math we need three sizes (text, script, scriptscript), so
    basically three copies with differently scaled
    dimensions.
\stopitem
\startitem
    We also can have increasing sizes.
\stopitem
\startitem
    And even extensibles built from snippets.
\stopitem
\startitem
    The number of fonts can grow large.
\stopitem
\stopitemize

\stoptitle

\starttitle[title=Evolution]

\startitemize
\startitem
    The \TEX\ engine is only interested in dimensions of glyphs.
\stopitem
\startitem
    We're talking width, height, depth, italic correction (not in \OPENTYPE),
    and expansion factors (hz).
\stopitem
\startitem
    In base mode it also like to know kerns and ligatures.
\stopitem
\startitem
    In math we have no design sizes so we can as well just scale. So we started
    doing just that
\stopitem
\startitem
    And then we could as do that for text. Welcome to compact fonts.
\stopitem
\stopitemize

\stoptitle

\starttitle[title=Evolution]

\startitemize
\startitem
    In text we have \typ {\glyphscale}, \typ {\glyphxscale} (extend), \typ
    {\glyphyscale} (squeeze).
\stopitem
\startitem
    We also have \typ {\glyphslant}, \typ {\glyphweight}
\stopitem
\startitem
    In math we have and  \typ {\Umathxscale} and \typ {\Umathyscale}.
\stopitem
\startitem
    We also have \typ {\glyphtextscale}, \typ {\glyphscriptscale} and
    \typ {\glyphscriptscriptscale}.
\stopitem
\startitem
    Text glyphs have expansion (and compression) properties for hz as well as
    protrusion specifications. For base mode we still have kerns and ligatures.
\stopitem
\stopitem
\startitem
    Math glyphs have a lot of extra properties: a \type {smaller}, \typ {mirror},
    \typ {flat accent}, \typ {next in size}, \typ {top anchor}, \typ {bottom anchor},
    \typ {top left}, \typ {top right}, \typ {bottom left}, \typ {bottom right}
    kerns, four (rather useless) staircase arrays, extensible recipes and italic
    correction, \typ {top}, \typ {bottom}, \typ {left} and \typ {right} margins,
    \typ {top} and \typ {bottom} overshoots, inner \type {x} and \typ {y} offsets.
\stopitem
\stopitemize

\stoptitle

\starttitle[title=Impact]

\startitemize
\startitem
    We need less definitions at the \TEX\ end because we scale on the fly.
\stopitem
\startitem
    We do still have instances for slant and weight variants because there
    the macro package decides impact on dimensions.
\stopitem
\startitem
    In math we resolve smaller (actually design) sizes on the fly too.
\stopitem
\startitem
    When processing \OPENTYPE\ features we already operated on multiple
    dimensions: instance (defaults), dynamic (smallcaps etc) and state (init,
    etc) features. We now also take scales, weight and slant into account.
\stopitem
\startitem
    Of course more calculations take place but in the end it gets compensated by
    less passes and memory usage.
\stopitem
\startitem
    We need different code paths at the macro level and might decide to remove
    the old method which then makes some more cleanup possible.
\stopitem
\stopitemize

\stoptitle

\starttitle[title=Tricky]

\startitemize
\startitem
    The backend has to work a bit harder to make sure that we don't mess up the
    \PDF\ width arrays and keep drift due to scaling in glyph streams
    within bounds.
\stopitem
\startitem
    Because we share font resources we need to handle font specific scales,
    slants, weights, as well as glyph bound scales, slant, weight and
    expansion.
\stopitem
\startitem
    We also have kerns (when encoded in glyphs), x- and y-offsets, virtual
    font (glyph) properties, etc. All this can interfere.
\stopitem
\startitem
    {\showglyphs Runtime {\glyphslant200 slant} and {\glyphweight200 weight} are okay for small values.}
\stopitem
\startitem
    But for instance spacing also matters so there is more involved.
\stopitem
\stopitemize

\stoptitle

\stopdocument

